\section{Anexos}

\subsection{Distribución espacial del PBI real per cápita relativo}

\begin{figure}[H]
\centering
\includegraphics[width=0.94\textwidth]{figura_A1_mapas_pbi.png}
\caption{PBI real per cápita relativo inicial y cambio acumulado}
\label{fig:anexo_mapas_pbi}
\begin{flushleft}
\footnotesize
\textit{Nota.} El primer panel representa el logaritmo del PBI real per cápita
relativo en 2007 y el segundo su cambio acumulado entre 2007 y 2019.
\end{flushleft}
\end{figure}

\subsection{Dispersión espacial del PBI per cápita}

\begin{figure}[H]
\centering
\includegraphics[width=0.90\textwidth]{figura_A2_moran_pbi.png}
\caption{Diagramas de dispersión de Moran para el PBI per cápita relativo}
\label{fig:anexo_moran_pbi}
\begin{flushleft}
\footnotesize
\textit{Nota.} Los diagramas corresponden al nivel inicial y al cambio
acumulado del indicador económico.
\end{flushleft}
\end{figure}

\subsection{Coeficientes locales de convergencia económica}

\begin{figure}[H]
\centering
\includegraphics[width=0.78\textwidth]{figura_5_gwr_beta_pbi.png}
\caption{Distribución territorial del coeficiente local del PBI per cápita}
\label{fig:anexo_beta_pbi}
\begin{flushleft}
\footnotesize
\textit{Nota.} Los valores negativos representan convergencia beta. Una mayor
magnitud absoluta indica una convergencia local más intensa.
\end{flushleft}
\end{figure}

\subsection{Pruebas de dependencia espacial}

\begin{table}[H]
\centering
\caption{Pruebas LM/RS de dependencia espacial}
\label{tab:anexo_lm}
\scriptsize
\begin{tabular}{lllrr}
\toprule
\textbf{Variable} &
\textbf{Matriz} &
\textbf{Prueba} &
\textbf{Estadístico} &
\textbf{$p$-valor} \\
\midrule
Cambio IDH & Banda inversa & RS error ajustado & 0.0329 & 0.8560 \\
Cambio IDH & Banda inversa & RS rezago ajustado & 0.0008 & 0.9770 \\
Cambio IDH & Queen & RS error ajustado & 2.3093 & 0.1286 \\
Cambio IDH & Queen & RS rezago ajustado & 1.8635 & 0.1722 \\
\midrule
Cambio $\ln$ PBIpc & Banda inversa & RS error ajustado & 4.1951 & 0.0405 \\
Cambio $\ln$ PBIpc & Banda inversa & RS rezago ajustado & 2.9907 & 0.0837 \\
Cambio $\ln$ PBIpc & Queen & RS error ajustado & 6.2672 & 0.0123 \\
Cambio $\ln$ PBIpc & Queen & RS rezago ajustado & 2.9723 & 0.0847 \\
\bottomrule
\end{tabular}
\begin{flushleft}
\footnotesize
\textit{Nota.} Las pruebas robustas permiten distinguir entre dependencia
espacial en el error y dependencia asociada con el rezago espacial. Para el
PBI per cápita, la evidencia se concentra principalmente en el componente de
error.
\end{flushleft}
\end{table}
